\documentclass{cv}

\usepackage[utf8]{inputenc}
\usepackage[ngerman]{babel}
\usepackage{hyperref}

\hypersetup{
	colorlinks=true,
	urlcolor=myblue,
	linkcolor=myblue,
	citecolor=myblue,
	pdftitle={Dr. André Kaderli -- Motivationsschreiben (Dozentur Quantencomputing)},
	pdfauthor={André Kaderli}
}

\renewcommand{\doclabel}{Motivationsschreiben}

% Rote Platzhalter \tbd{...} vor dem Versand ausfuellen oder entfernen.

\begin{document}

	\name{Dr. André Kaderli}
	\begin{contactline}
		Obere Bahnhofstrasse 22M, 3700 Spiez\sep
		\href{mailto:andre.kaderli@gmail.com}{andre.kaderli@gmail.com}\sep
		\href{tel:+41793323931}{+41\,79 332 39 31}\sep
		\href{https://www.linkedin.com/in/andre-kaderli/}{linkedin.com/in/andre-kaderli}
	\end{contactline}

	\recipient{%
		ZHAW Zürcher Hochschule für Angewandte Wissenschaften\\
		School of Engineering\\
		Institut für Angewandte Mathematik und Physik (IAMP)\\
		\tbd{Ansprechperson}\\
		Technikumstrasse 9, 8401 Winterthur \tbd{prüfen}}

	\placedate{Spiez, \tbd{Datum}}

	\lettersubject{Bewerbung für die Dozierendenstelle Quantencomputing und Quantum Machine Learning}

	\begin{letterbody}
	Sehr geehrte\tbd{r Frau/Herr Name}

	Vielen Dank für Ihre Rückmeldung auf meine Spontanbewerbung und den Hinweis auf die Dozierendenstelle am IAMP. Gerne nutze ich diese Chance: Die Verbindung von Forschung in Quantencomputing und Quantum Machine Learning mit der Lehre in Mathematik entspricht genau dem, was mich fachlich antreibt.

	Mein wissenschaftliches Fundament ist die Quantentheorie. An der ETH Zürich habe ich theoretische und mathematische Physik studiert (MSc mit Note 6.0) und dabei auch Quanteninformation und Quantencomputing kennengelernt (\tbd{Kursname}). Anschliessend habe ich am Max-Planck-Institut für Gravitationsphysik und an der Humboldt-Universität zu Berlin mit summa cum laude promoviert. Dort habe ich mit analytischen und numerischen Methoden sowie rekursiven Algorithmen an Problemen der Quantengravitation gearbeitet. Lineare Algebra in Hilberträumen, Gruppentheorie und Wahrscheinlichkeitstheorie, also die mathematische Sprache des Quantencomputings, sind für mich vertrautes Werkzeug.

	Seit über fünf Jahren setze ich Machine Learning, Statistik und mathematische Modellierung in Projekten für öffentliche Verwaltung und Privatwirtschaft um, heute als Partner bei der Novalytica AG. Dabei habe ich gelernt, reale Fragestellungen in quantitative Modelle zu übersetzen und Lösungen bis in den produktiven Betrieb zu bringen. Diese Erfahrung möchte ich in die anwendungsorientierte Forschung am IAMP einbringen. Mich interessiert vor allem, wo Quantum Machine Learning gegenüber klassischen Verfahren einen messbaren Vorteil bringt. Dazu möchte ich variationale Quantenalgorithmen und Quantum Kernels an realen Datensätzen untersuchen, systematisch mit klassischen Methoden vergleichen und gemeinsam mit Praxispartnern Projekte aufbauen, etwa mit Innosuisse.

	Lehren begleitet mich seit über 15 Jahren: von Übungsgruppen in Linearer Algebra an der ETH über Vorlesungsunterstützung in Funktionentheorie an der HU Berlin bis zu Weiterbildungskursen in Statistik und Datenanalyse für Fachleute. Mein Anspruch ist, mathematische Strenge mit anschaulichen Beispielen aus der Praxis zu verbinden, damit Studierende verstehen, wozu sie Analysis, Lineare Algebra, Statistik oder Numerik brauchen. Genau das verstehe ich unter Lehre an einer Hochschule für Angewandte Wissenschaften.

	Über die Gelegenheit, Ihnen meine Ideen für Forschung und Lehre in einem persönlichen Gespräch vorzustellen, freue ich mich sehr. \tbd{Verfügbarkeit: ab wann, gewünschtes Pensum}
	\end{letterbody}

	\closing{Freundliche Grüsse}{Dr. André Kaderli}

	\enclosures{Lebenslauf, Diplome, Arbeitszeugnisse \tbd{prüfen}}

\end{document}
